% =============================================================================
% Calibration & Fidelity Analysis Section
% Auto-generated by run_fidelity_analysis.py — edit as needed
% =============================================================================

\subsection{Calibration Against Real Sensor Traces}
\label{sec:fidelity}

To quantify the realism of STGen's sensor models, we benchmark
synthetic traces against the Intel Berkeley Research Lab
dataset~\cite{intel-berkeley-lab}, which contains approximately 2.3
million readings from 54 Mica2Dot sensors collecting temperature,
humidity, light, and voltage data at $\sim$31-second intervals over 37
days. After filtering implausible readings (e.g., negative humidity,
voltages outside 1.5--3.5\,V), 1,826,223 valid samples remain.

We compare four classes of metrics:
(i)~marginal distributional match (Kolmogorov--Smirnov $D$-statistic and
Jensen--Shannon divergence),
(ii)~temporal autocorrelation (ACF at lags $\ell \in \{1,5,10,20,50\}$),
(iii)~cross-modal coupling (temperature--humidity Pearson~$r$), and
(iv)~inter-arrival time distribution.
For each modality, we generate 100,000 synthetic samples using
STGen's physics-based models (Section~\ref{sec:sensor-models}), then
repeat with parameters \emph{fitted} to the Intel traces to show the
calibration ceiling.

%% -------------------------------------------------------------------
\subsubsection{Why Variation Matters: Inter- vs.\ Intra-Mote Decomposition}

A variance decomposition of the Intel data reveals that \textbf{95.3\%
of temperature variance, 92.3\% of humidity variance, and 87.9\% of
light variance is within-mote} (temporal), not between-mote (spatial).
The per-mote standard deviation averages 3.82\textdegree C for
temperature and 5.86\% for humidity; the between-mote standard
deviation of means is only 1.01\textdegree C and 1.76\%,
respectively.  This finding has two direct consequences for
traffic generation:

\begin{enumerate}[nosep]
  \item \textbf{Temporal persistence matters more than spatial diversity.}
    An i.i.d.\ noise model with the correct marginal moments will produce
    near-zero autocorrelation, destroying the serial structure that
    dominates real traces.
  \item \textbf{Regime switching is essential.}  Individual motes exhibit
    heavy-tailed temperature distributions (per-mote skewness up to
    $6.0$, kurtosis up to $18.5$), which cannot arise from a single
    stationary OU process.  These tails are caused by physical
    events---HVAC cycling, occupancy changes, and direct sunlight---that
    shift the mote's local mean.
\end{enumerate}

%% -------------------------------------------------------------------
\subsubsection{Model Architecture}

Based on this analysis, STGen uses a \emph{multivariate OU
architecture} with regime switching.  All four sensor modalities are
generated by coupled Ornstein--Uhlenbeck processes:

\paragraph{Temperature.}
Each mote~$i$ follows an OU process
$dT_i = \theta(\mu_i - T_i)\,dt + \sigma\,dW$,
with per-mote mean $\mu_i \sim \mathcal{N}(22.2, 1.0)$ and
regime-switching: at each step there is a 0.1\% probability that
$\mu_i$ shifts by $\Delta\mu \sim \mathcal{N}(0, 3)$, modelling HVAC
or sunlight events.  Calibrated parameters: $\theta = 0.002$,
$\sigma = 0.15$, rate cap $2$\textdegree C/min.

\paragraph{Humidity.}
An OU process mean-reverts toward a temperature-coupled target
$h^* = 53.8 - 0.64\,T$ with $\theta_h = 0.003$ and $\sigma_h = 0.2$,
providing both cross-correlation with temperature and strong temporal
autocorrelation.

\paragraph{Light.}
A bimodal regime-switching OU model alternates between a \emph{dark}
attractor ($\mu = 28$\,lux) and a \emph{bright} attractor
($\mu = 639$\,lux), with transition probabilities 2\% and 1\% per step
respectively, matching the Intel data where 37.5\% of readings fall
below 100\,lux.

\paragraph{Voltage.}
An OU process ($\theta_v = 0.001$, $\sigma_v = 0.005$) with slow
deterministic discharge drift ($-10^{-5}$\,V/step), centred at
$\mu = 2.56$\,V to match Mica2Dot Li-ion cells.

%% -------------------------------------------------------------------
\subsubsection{Marginal Distribution Fidelity}

Table~\ref{tab:fidelity-marginal} summarises the marginal statistics.
STGen's OU temperature model achieves mean agreement within 0.4\%
($22.01$\textdegree C vs.\ $22.43$\textdegree C) with KS~$D = 0.142$.
Voltage is the closest match overall (KS~$D = 0.085$, JSD~$= 0.017$),
as the OU drift model closely replicates the narrow Intel voltage
distribution ($2.56 \pm 0.12$\,V).

Humidity achieves KS~$D = 0.204$ with the correct mean ($39.4$\%
vs.\ $39.5$\%), though variance is still under-estimated
($\sigma = 2.98$ vs.\ $6.28$) because the OU noise coupling is
conservative.  Light shows the largest gap (KS~$D = 0.484$)
because the bimodal OU within-regime variance is smaller than the
real data's heavy tails ($\sigma = 541$ lux).

\input{tables/fidelity_marginal}

%% -------------------------------------------------------------------
\subsubsection{Temporal Autocorrelation}

Table~\ref{tab:fidelity-acf} shows that all four modalities now
exhibit strong temporal persistence, closely matching the Intel traces.
Temperature ACF at lag~50 is 0.910 vs.\ 0.876 (Intel), light reaches
0.826 vs.\ 0.913, and voltage 0.832 vs.\ 0.992.  The key improvement
is in humidity and light, which previously had ACF~$\approx 0$ at all
lags under i.i.d.\ models.  The OU architecture produces lag-1
autocorrelations exceeding 0.993 for all modalities (Intel: $>$0.998),
confirming that the serial structure of real sensor readings is
faithfully replicated.

\input{tables/fidelity_acf}

%% -------------------------------------------------------------------
\subsubsection{Cross-Correlation and Inter-Arrival Times}

The Intel data shows a moderate negative temperature--humidity coupling
($r = -0.43$).  STGen's OU-coupled model produces $r = -0.35$,
capturing the sign and approximate magnitude of the anti-correlation.
The remaining gap arises because the real relationship is non-linear
and influenced by diurnal and occupancy cycles not modelled by a
first-order linear coupling.

The nominal Intel sampling interval is 31\,s, but real inter-arrival
times exhibit substantial jitter ($\mu = 47.5$\,s, $\sigma = 35.5$\,s)
from retransmissions, clock drift, and packet loss---consistent with
low-power wireless behaviour.  STGen's Weibull IAT model (with
$k = 0.8$, scale matching the 31\,s base) produces $\mu = 35.4$\,s,
$\sigma = 44.6$\,s, capturing the right-skewed, bursty structure of
real IATs more faithfully than a fixed-interval model.

%% -------------------------------------------------------------------
\subsubsection{Overall Fidelity}

The composite fidelity score---averaging the KS, JSD, and ACF
sub-scores---reaches $\mathbf{0.865}$ out of 1.000 across all four
modalities.  Temperature and voltage score highest (0.925 and 0.947),
while light is lowest (0.734) due to its extreme bimodal variance.
With parameters fitted directly to Intel traces (Phase~2 calibration),
the score remains comparable at 0.860, indicating that the
\emph{default} model architecture already captures the dominant
structure of real sensor data; further marginal gains require matching
the exact distributional shape (heavy tails, non-Gaussian modes).

\input{tables/fidelity_improvement}

%% -------------------------------------------------------------------
\subsubsection{Remaining Gap: Temperature Variance}

The primary remaining gap is temperature standard deviation
($1.65$\textdegree C vs.\ $4.19$\textdegree C).  The Intel data's
heavy tails (kurtosis~$18.5$) arise from physical events---HVAC
cycling, occupancy-driven heating, direct sunlight exposure---that
produce transient excursions far from the mote's baseline.  While our
regime-switching mechanism (0.1\% shift probability) introduces some
tail behaviour, the $2$\textdegree C/min rate cap still constrains the
maximum excursion within any OU trajectory.  Relaxing this cap or
adding explicit event-driven spikes (modelled as Poisson arrivals with
random amplitude) would further close the gap.
